\subsection{绝对半单模的刻画}

命题 4. —— 设 K 是一个交换域，A 是一个 K-代数。

\begin{enumerate}
\def\labelenumi{\alph{enumi})}
\item
  设 M 是一个半单 A-模。A-模 M 绝对半单当且仅当属于 M 的支集的每个单模都是绝对半单的。
\item
  设 S 是一个单 A-模，D 是它的交换子。下列性质等价：
\end{enumerate}

\begin{enumerate}
\def\labelenumi{(\roman{enumi})}
\item
  A-模 S 是绝对半单的。
\item
  K-代数 \(\mathbf{D}\) 是绝对半单的。
\item
  K-代数 D 是一个体，在 K 上的次数有限，且其中心是 K 的一个可分扩张。
\end{enumerate}

断言 a) 由 VIII, 第 229 页的命题 1, a) 得出。设 S 与 D 如 b) 中所述，并设 L 是 K 的一个扩张。\(A_{(L)}\) -模 \(S_{(L)}\) 是半单的当且仅当环 \(D_{(L)}\) 是半单的（VIII, 第 222 页，命题 8, c)）。这就证明了 (i) 与 (ii) 的等价性，而 (ii) 与 (iii) 的等价性由定理 1 得出，因为 D 是一个体。

推论. —— 设 K 是一个交换域，\(A_{1}\) 与 \(A_{2}\) 是两个 K-代数，\(M_{1}\) 是一个绝对半单 \(A_{1}\) -模，而 \(M_{2}\) 是一个半单 \(A_{2}\) -模。则 \(M_{1} \otimes_{K} M_{2}\) 是环 \(A_{1} \otimes_{K} A_{2}\) 上的一个半单模。

模 \(M_{1}\) 是绝对半单的单 \(A_{1}\) -模的直和（命题 4）。因此只需在模 \(M_{1}\) 与 \(M_{2}\) 都是单的情形证明断言。用 \(D_{1}\) 与 \(D_{2}\) 表示它们的交换子。K-代数 \(D_{1}\) 是绝对半单的（同前所引）；由 VIII, 第 234 页的推论 1，K-代数 \(D_{1} \otimes_{K} D_{2}\) 是半单的。于是由 VIII, 第 215 页的推论 1 得出，\((\mathrm{A}_{1} \otimes_{\mathrm{K}} \mathrm{A}_{2})\) -模 \(M_{1} \otimes_{K} M_{2}\) 是半单的。

